% prologo.tex — versión con notación coherente y conexión actuarial (pdfLaTeX)
\chapter*{Prólogo}
\addcontentsline{toc}{chapter}{Prólogo}

\begingroup
\setlength{\parindent}{0pt}
\setlength{\parskip}{6pt}
\emergencystretch=3em
\sloppy

\begin{flushright}
	\emph{\guillemotleft Tu única brújula fiable es la curvatura de tu propia incertidumbre.\guillemotright}
\end{flushright}

Este libro parte de una tensión de todos los días: el mundo sigue una distribución verdadera $p(x)$ y nosotros trabajamos con un modelo $p_\theta(x)$ con parámetro $\theta$.

Bayes da una receta de actualización limpia usando el modelo:
\[
\pi(\theta\mid x)\;\propto\;\pi(\theta)\,p_\theta(x).
\]
Entonces, ¿para qué geometría de la información (GI)? Para ponerle \emph{geometría} al espacio de distribuciones y convertir varias reglas sueltas en \textbf{principios invariantes}: rectas, proyecciones y sensibilidades que no dependen de cómo parametrices el modelo.

\paragraph{¿Qué es reparametrizar, en la práctica, y para qué sirve?}
Cambiar de coordenadas para trabajar mejor. Tres casos comunes:
\begin{itemize}
	\item \textbf{Logístico/GLM:} si una probabilidad $\pi\in(0,1)$, usar el logit $\eta=\log(\pi/(1-\pi))=x^{\top}\beta$ elimina límites y estabiliza la optimización.
	\item \textbf{Normal:} usar \emph{precisión} $\tau=1/\sigma^2$ en lugar de varianza positiva; cómodo en Bayes y mejor condicionado en optimización.
	\item \textbf{Positividad:} transformar $\lambda>0$ a $u=\log\lambda$ (Poisson, tasas, escalas) para optimizar en $u$ sin violar restricciones.
\end{itemize}
Estos cambios no deberían alterar \emph{las conclusiones} cuando no cambian la información disponible. GI ofrece herramientas para exigir esa coherencia entre coordenadas.

\textbf{Qué añade GI y por qué importa}

\begin{enumerate}
	\item \textbf{Métrica elegida por invariancia (Chentsov).} Si exiges que las distancias no cambien al aplicar transformaciones estocásticas que no pierden información (estadísticos suficientes), la métrica Riemanniana queda determinada (salvo constante): \emph{Fisher--Rao}. Representante típico:
	\[
	g_{ij}(\theta)\;=\;\mathrm{E}_{\theta}\!\big[\partial_i \log p_\theta(X)\,\partial_j \log p_\theta(X)\big].
	\]
	No es cuestión de gusto: es una \emph{unicidad} por invariancia. La sensibilidad refleja información, no coordenadas.
	
	\item \textbf{Dos “rectas” canónicas (dualidad e/m).} En familias exponenciales, el espacio es dualistamente plano: hay dos líneas rectas naturales. Una preserva promedios (coordenadas de expectativa); la otra, parámetros naturales. Explica por qué existen \textbf{dos maneras legítimas de actualizar} y cuándo conviene cada una.
	
	\item \textbf{Proyecciones KL con “pitagórico”.} Al aproximar $p$ con una familia $\mathcal{Q}$, GI propone \emph{proyectar} por KL. Se obtiene
	\[
	D_{\mathrm{KL}}(p\Vert q)\;=\;D_{\mathrm{KL}}(p\Vert q^\star)\,+\,D_{\mathrm{KL}}(q^\star\Vert q),
	\]
	donde $q^\star$ es la proyección de $p$ en $\mathcal{Q}$. Da dirección de mejora y criterio claro de convergencia.
	
	\item \textbf{Gradiente que respeta la información (natural).} Con la métrica de Fisher:
	\[
	\dot{\theta}\;=\;-\,I(\theta)^{-1}\,\nabla_{\!\theta} L(\theta),
	\]
	obtienes pasos invariantes a reparametrizaciones, con menos zigzag y mejor alineados con la información.
	
	\item \textbf{Puentes con herramientas que ya usas.}
	\begin{itemize}
		\item \textbf{AIC} como \emph{riesgo KL esperado}: elegir menor AIC aproxima al modelo que deja más cerca (en KL) de $p$ a futuro.
		\item \textbf{Prior de Jeffreys} $\pi_J(\theta)\propto \sqrt{\det I(\theta)}$: nace de pedir \emph{invariancia} (volumen Riemanniano), en el mismo espíritu de Chentsov.
		\item \textbf{EM y VI} como \emph{proyecciones alternadas}: pasos I/M en variedades afines; la monotonicidad viene del “pitagórico” de KL.
	\end{itemize}
\end{enumerate}

\textbf{Medir el desajuste entre verdad y modelo}

Para cuantificar el \emph{desajuste} entre $p$ y $q$ usamos divergencias. La más común es la KL:
\[
D_{\mathrm{KL}}(p\Vert q)\;=\;\sum_x p(x)\,\log\frac{p(x)}{q(x)}\,.
\]
Una \emph{divergencia} $D(p\Vert q)$ no es (necesariamente) simétrica: los errores tienen dirección; no cuesta igual confundir $p$ con $q$ que confundir $q$ con $p$\footnote{O, como diría un boricua: \emph{es peor que fafalte a que sosobre}.}. KL mide el \textbf{costo del desajuste} al describir $p$ como si el mundo fuese $q$. Esa asimetría te permite penalizar lo que realmente te importa.

\textbf{Dos esquemas de actualización (con lectura actuarial)}

Imagina que, tras ver datos nuevos, quieres acercar $p_\theta(x)$ a la distribución verdadera. Hay dos esquemas muy usados, con interpretación directa en \emph{credibilidad}:

\begin{enumerate}
	\item \textbf{Esquema por momentos (credibilidad clásica, B{\"u}hlmann--Straub).}
	Ajustas el modelo para que respete ciertos \textbf{promedios} o totales confiables (tasa global, proporción, total esperado). El estimador típico combina experiencia propia y base colectiva:
	\[
	\hat{\theta}\;=\;Z\,\bar{X}\;+\;(1-Z)\,m,\qquad Z\in[0,1].
	\]
	En GI, corresponde a una \emph{e-proyección} (preserva momentos).
	
	\item \textbf{Esquema conservador (credibilidad bayesiana).}
	Das un paso que cambia lo necesario según la evidencia, manteniendo razonable la incertidumbre global. Con conjugadas:
	\[
	\mathrm{E}[\theta\mid \mathrm{datos}]\;=\;\frac{n}{\,n+k\,}\,\bar{X}\;+\;\frac{k}{\,n+k\,}\,m,
	\]
	con $n$ tamaño muestral y $k$ la “fuerza” previa. Pocos datos $\Rightarrow$ más peso a $m$; muchos datos $\Rightarrow$ manda $\bar{X}$. En GI, se lee como \emph{m-proyección} (recta en parámetros naturales).
\end{enumerate}

No compiten: son \emph{herramientas complementarias}. La mirada geométrica te dice qué conserva cada una y por qué su combinación (como en EM/VI) funciona.

\textbf{Coherencia que no negocia}

Si reexpresas los datos sin perder información (estadísticos suficientes, “coarse-graining” honesto), tus conclusiones no deberían cambiar. Esa exigencia de invariancia selecciona Fisher--Rao y hace que los métodos derivados sean realmente independientes de la etiqueta de las coordenadas.

\textbf{Qué te llevas de este libro}

\begin{itemize}
	\item Un \textbf{idioma geométrico} donde “rectas”, “proyecciones” y “curvaturas” sobre distribuciones explican y mejoran prácticas que ya conoces.
	\item \textbf{Criterios invariantes} (Chentsov, Fisher--Rao) para no cambiar de respuesta por reparametrizar.
	\item \textbf{Proyecciones con garantía} (pitagórico de KL) para aproximar, alternar pasos y saber por qué converges.
	\item \textbf{Pasos que respetan la información} (gradiente natural) y conexiones claras con AIC, Jeffreys, EM y VI.
\end{itemize}

No buscamos efectos especiales. Buscamos \emph{orden}: medir, comparar y moverte en el espacio de modelos de manera coherente, independientemente de cómo nombres las cosas. Si una idea sobrevive a cambiarle las coordenadas, probablemente vale la pena.

\begin{flushright}
	\emph{Que tu gradiente sea suave y tu sorpresa, cada vez menor.}
\end{flushright}

\endgroup
